%%%%%%%%%%%%%%
%
%  第一部分 软件概述
%
%%%%%%%%%%%%%%
\section{软件概述}

\subsection{项目背景}
\par
超声检测通过分析材料内部反射回波的传播时间、幅值和频率特征，判断工件中是否存在裂纹、夹渣、未熔合、气孔或分层等缺陷。A-Scan 以时间为横轴、回波幅值为纵轴，是超声检测中最直接、最常用的显示形式。现场检测时，操作人员需要根据工件材料和几何形状选择探头与耦合方式，结合声速、声程、增益、闸门和游标完成缺陷定位与定量，并保留检测过程、缺陷记录和结论依据。
\par
“\SoftwareName（\Version）”围绕上述工作流程，将波形采集、参数配置、闸门判定、游标测量、校准、缺陷记录、数据回放、统计分析和报告输出集中在同一桌面环境中。软件采用深色工业界面与绿色网格波形显示，支持多闸门、TCG 补偿、B/C-Scan 成像、FFT 频谱分析和 SQLite 会话存储，为教学演示、工艺试验、算法验证和检测数据管理提供统一操作平台。

\subsection{编写目的}
\par
本说明书用于说明软件的安装、登录、参数设置、校准、检测、记录、回放和报告导出方法，并将常用超声检测作业步骤与软件操作对应起来。阅读本说明书后，操作人员应能够根据工件与探头条件完成声速和零点校准，设置适合的声程与增益，使用 Gate A/B/C 和游标定位缺陷，保存可追溯的检测会话，并输出包含检测参数、波形和缺陷信息的 HTML 或 PDF 报告。

\subsection{适用对象与使用范围}
\par
本软件面向无损检测人员、质量工程师、检测设备与算法研发人员、高校教师和科研人员。实际检测任务中，操作人员应依据工件材料、焊接工艺、验收标准和检测工艺卡确定探头频率、晶片尺寸、入射角度、耦合方式、扫查范围、扫查速度及验收等级；软件负责将这些条件转化为可执行、可复现的采集与记录流程。
\par
软件的适用范围如\cref{tab:software_scope}所示。各功能既可以独立使用，也可以在典型检测流程中连续执行。

\begin{table}[htbp]
	\centering
	\caption{软件适用范围}
	\label{tab:software_scope}
	\begin{tabular}{p{0.27\textwidth}p{0.61\textwidth}}
		\toprule
		\textbf{功能类别} & \textbf{主要用途} \\
		\midrule
		波形采集与显示 & A-Scan 实时显示、冻结、单次采样、峰值保持、参考波形比较和自动缩放。 \\
		检测参数与校准 & 增益、声程、探头频率、刷新间隔、声速、零点延迟及校准文件管理。 \\
		闸门与测量 & Gate A/B/C 门限判定、峰值捕捉、深度与水平距离计算、A/B 游标差值测量。 \\
		信号分析 & FFT 频谱、窗函数、带通/低通/高通滤波、信号平均、特征提取和参考波形对比。 \\
		缺陷分析 & 缺陷深度与波幅统计、DAC 曲线、TCG 补偿和超限判定。 \\
		扫查成像 & B-Scan 时间-位置显示、C-Scan 平面位置显示和光栅扫描坐标管理。 \\
		数据与报告 & SQLite 会话保存、波形分块、逐帧回放、JSON/CSV 导出、HTML/PDF 报告。 \\
		权限与追溯 & 本地用户登录、四类角色、操作审计、报告签名和文件哈希。 \\
		\bottomrule
	\end{tabular}
\end{table}

\subsection{软件特点}
\par
软件采用 MVC 架构，模型层负责参数、信号仿真、闸门、TCG、统计和持久化，视图层负责波形与面板显示，控制器负责连接用户操作和业务逻辑。该结构使波形算法、界面控件和数据存储可以分别维护，便于后续接入不同采集设备或扩展检测标准。
\par
波形区域使用 PyQtGraph 绘制原始信号、包络、参考波形和峰值保持曲线，闸门线、门限线和游标均可直接拖动。界面在窗口大小变化时保持坐标范围稳定，检测表格、频谱面板、统计面板、回放面板和扫描成像面板可从“视图”菜单按需显示。
\par
检测会话以 SQLite 保存参数、缺陷、报告和波形分块索引，原始波形按会话写入压缩 NPZ 文件。报告采用同一 HTML 内容输出为网页或 PDF，并记录报告文件 SHA-256。审计事件使用前序哈希链连接，用户、角色、操作时间和报告签名均可追溯。

\subsection{运行基础与术语}
\par
软件支持 Windows 10/11、macOS 和 Linux，运行环境以 Python 3.10 及以上版本为基础。检测会话、用户、缺陷和报告元数据保存在 SQLite 数据库中，波形数组按会话写入压缩 NPZ 文件。为保证长时间采集、较高采样率和 B/C-Scan 显示稳定，推荐使用 8 GB 以上内存、10 GB 以上可用硬盘空间以及 1366$\times$768 以上显示分辨率。
\par
软件使用的运行配置、检测参数、测量术语和分析概念如\cref{tab:runtime_terms}所示。表中的检测参数直接对应右侧控制面板，术语部分用于统一后续操作说明和报告数据的含义。

\begin{longtable}{p{0.18\textwidth}p{0.27\textwidth}p{0.43\textwidth}}
	\caption{运行基础、检测参数与术语}
	\label{tab:runtime_terms}\\
	\toprule
	\textbf{类别} & \textbf{项目或术语} & \textbf{说明} \\
	\midrule
	\endfirsthead

	\toprule
	\textbf{类别} & \textbf{项目或术语} & \textbf{说明} \\
	\midrule
	\endhead

	\bottomrule
	\endlastfoot

	\multicolumn{3}{l}{\textbf{运行环境}} \\
	\midrule
	操作系统 & Windows 10/11、macOS、Linux & 图形界面基于 Qt 5，程序代码保持跨平台兼容。 \\
	Python & 3.10 及以上 & 推荐使用独立虚拟环境安装运行依赖，便于升级和回退。 \\
	内存 & 8 GB 及以上 & 长会话、较高采样率和 B/C-Scan 显示建议使用 16 GB。 \\
	硬盘 & 10 GB 可用空间 & 波形分块和报告文件随检测会话数量增长。 \\
	显示分辨率 & 1366$\times$768 及以上 & 1920$\times$1080 可同时显示更多面板内容。 \\
	Python 依赖 & PyQt5、PyQtGraph、NumPy、SciPy & 版本要求见项目目录中的 \texttt{requirements.txt}。 \\
	数据存储 & SQLite 与本地文件系统 & 数据库保存索引和元数据，NPZ 文件保存波形数组。 \\
	\midrule

	\multicolumn{3}{l}{\textbf{检测参数与测量术语}} \\
	\midrule
	A-Scan & A 型显示，以时间为横轴、幅值为纵轴的单通道波形。 & 主波形区显示原始信号与包络。 \\
	UT & Ultrasonic Testing，超声检测。 & 软件面向的检测领域。 \\
	增益 & 接收信号的整体放大倍数，单位为 dB。 & 控制回波幅值，不改变传播时间。 \\
	声程 & 波形窗口对应的最大检测距离，单位为 mm。 & 与声速共同决定时间轴长度。 \\
	声速 & 超声波在材料中的传播速度，单位为 m/s。 & 用于时间与深度换算。 \\
	零点延迟 & 探头、楔块和电路引入的固定时间偏移，单位为 $\mu$s。 & 校准浅表缺陷读数。 \\
	闸门 & 在指定时间或深度范围内检测回波的窗口。 & Gate A/B/C 分别设置起点、宽度和门限。 \\
	门限 & 触发记录或判定的最小包络幅值。 & 与闸门内峰值比较。 \\
	游标 & 用于读取任意位置时间、深度和幅值的测量线。 & A/B 两条游标显示差值与距离。 \\
	\midrule

	\multicolumn{3}{l}{\textbf{分析与成像术语}} \\
	\midrule
	FFT & Fast Fourier Transform，快速傅里叶变换。 & 将时域波形转换为频率谱。 \\
	DAC & Distance-Amplitude Correction，距离-波幅校正。 & 生成深度-幅值参考曲线。 \\
	TCG & Time-Corrected Gain，时间增益补偿。 & 按深度分段补偿回波幅值。 \\
	B-Scan & 以扫描位置为纵轴、时间为横轴的二维波形堆叠图。 & 显示检测过程的时间变化。 \\
	C-Scan & 以 X/Y 平面位置为坐标的幅值分布图。 & 显示光栅扫描区域的缺陷分布。 \\
\end{longtable}
